\chapter{旅游产品体系刻画与耦合分析}\label{ux65c5ux6e38ux4ea7ux54c1ux4f53ux7cfbux523bux753bux4e0eux8026ux5408ux5206ux6790}

\section{旅游产品体系的结构}\label{ux65c5ux6e38ux4ea7ux54c1ux4f53ux7cfbux7684ux7ed3ux6784}

\subsection{类型构成}\label{ux7c7bux578bux6784ux6210}

基于清洗后的 13,512 条 POI，南海区旅游产品可归入 11 个类型（表 5.1）。其中，``公园绿地''和``自然景观''合计占 70.8\%，说明当前旅游供给首先表现为绿色开放空间和自然游憩空间的集中；而``非遗体验''\,``文化场馆''和``特色街区''合计仅占 3.8\%，反映出可直接转化为文化体验的产品供给仍然偏少。``其他''类占 14.9\%，主要作为分类口径中的补充项，本文在后续解释中不对其作过度文化化判断。


\begin{table}[H]
\centering
\msttableentry{5.1}{13,512 条 POI 按 11 大类的构成}
\noindent\textbf{表 5.1　13,512 条 POI 按 11 大类的构成}\par\vspace{0.35em}
\begin{center}
{\scriptsize
\setlength{\tabcolsep}{2.5pt}
\renewcommand{\arraystretch}{1.12}
\resizebox{\textwidth}{!}{%
\begin{tabular}{@{}lrrrrrrrrrrrr@{}}
\toprule
\textbf{指标} & \textbf{公园绿地} & \textbf{自然景观} & \textbf{其他} & \textbf{宗教场所} & \textbf{人文古迹} & \textbf{文化场馆} & \textbf{休闲娱乐} & \textbf{体育设施} & \textbf{非遗体验} & \textbf{教育研学} & \textbf{特色街区} & \textbf{合计} \\
\midrule
数量 & 5,441 & 4,123 & 2,018 & 481 & 420 & 318 & 242 & 175 & 137 & 104 & 53 & 13,512 \\
占比 & 40.3\% & 30.5\% & 14.9\% & 3.6\% & 3.1\% & 2.4\% & 1.8\% & 1.3\% & 1.0\% & 0.8\% & 0.4\% & 100\% \\
\bottomrule
\end{tabular}
}
}
\end{center}
\end{table}

\begin{figure}
\centering
\includegraphics[width=5.86042in,height=3.25625in]{ms_thesis/image5.png}
\caption{13,512 条 POI 的类型构成与镇街分布}
\end{figure}

\subsection{镇街分布}\label{ux9547ux8857ux5206ux5e03}

按镇街统计，桂城街道（4,174）、里水镇（2,625）、狮山镇（2,233）、大沥镇（2,080）四个东北部镇街合计占据 82.1\% 的 POI 供给；西南部的丹灶镇（1,086）、西樵镇（1,055）、九江镇（259）则在供给侧明显偏低。九江镇 POI 数量占全区 1.92\%，是 POI 供给最为稀薄的镇街。这一结构为第 6 章关于东部旅游供给较强、西南部文化载体较密的空间错位讨论提供基础。

\section{文化---旅游耦合桥梁：物质遗产为主、非遗为补充}\label{ux6587ux5316ux65c5ux6e38ux8026ux5408ux6865ux6881ux7269ux8d28ux9057ux4ea7ux4e3aux4e3bux975eux9057ux4e3aux8865ux5145}

\subsection{桥梁的操作化设计}\label{ux6865ux6881ux7684ux64cdux4f5cux5316ux8bbeux8ba1}

本研究将文化记忆与旅游供给区分为两条分析线：知识图谱主要用于表征文化记忆侧的语义联系，POI 与评论数据主要用于表征旅游供给侧的现实热度。二者之间不能直接等同，因此需要设置一组可被复核的``耦合桥梁''。在本研究中，桥梁指具有明确名称、较稳定空间指向，并能够同时进入文化叙事和旅游空间分析的文化载体。

桥梁的核心样本由 165 条文化载体组成，其类型构成见表 3.3。选择这些载体作为主样本，主要基于三点考虑：其一，多数载体具有较明确的空间位置，能够与 POI 及后续空间网格进行对应；其二，文物保护单位、历史文化名村、风景名胜区以及非遗项目等载体具有不同层级的官方认定，可作为官方认证指数的依据；其三，载体名称相对稳定，有利于与知识图谱中的实体进行名称匹配。

\subsection{动态补充：依托空间的非遗项目}\label{ux52a8ux6001ux8865ux5145ux4f9dux6258ux7a7aux95f4ux7684ux975eux9057ux9879ux76ee}

90 项非遗项目中，有 36 项能够对应到较明确的空间载体，例如九江双蒸博物馆、西樵醒狮训练基地和松塘龙舟训练场等。这一部分非遗项目并入 165 条核心载体，参与空间匹配与旅游热度计算。其余 54 项主要依托传承人登记地或镇街信息，空间边界较弱，因此不直接与单个 POI 作一一对应，而是作为文化记忆侧的补充信息进入指数计算和后续规划判断。

\subsection{两条数据线的汇流}\label{ux4e24ux6761ux6570ux636eux7ebfux7684ux6c47ux6d41}

从管理条线看，文旅系统中的文物保护单位、遗址公园和非遗项目，住建系统中的风景名胜区、历史文化名镇名村与传统村落，虽然来源不同，但都可以在地方文化记忆与旅游空间之间形成连接。本文将其统一整理为文化载体样本，并不是把所有资源简单并列，而是将其放入同一套``文化记忆---官方认证---旅游热度''的比较框架中。这样可以减少单一依赖非遗或单一依赖景点名单带来的偏差，也更符合南海区文化资源多条线并存的实际情况。

\subsection{匹配规则}\label{ux5339ux914dux89c4ux5219}

匹配过程分为两个方向。第一，向文化记忆侧匹配，即以载体名称为入口，在知识图谱中检索相关实体，并结合名称简化和实体类型约束提高准确性。例如，``九江双蒸酒酿制技艺''可通过``双蒸酒''等核心词汇与图谱实体建立联系。第二，向旅游供给侧匹配，即以载体位置为中心，结合名称匹配与 500 m 缓冲范围内的 POI 进行对应。前者用于衡量文化记忆的可见度，后者用于衡量旅游供给的覆盖度。对于无法明确空间定位的项目，本文只将其纳入文化记忆侧，不将其强行转化为 POI 周边热度。

\section{文化---旅游量化指数与错位指数}\label{ux6587ux5316ux65c5ux6e38ux91cfux5316ux6307ux6570ux4e0eux9519ux4f4dux6307ux6570}

本节构造文化记忆指数、官方认证指数、旅游热度指数和错位指数四个指标，用于替代单纯的定性判断。需要说明的是，这些指标并不试图完整评价某一文化载体的价值，也不直接测度游客对文化内容的识别深度，而是服务于本研究的核心问题：识别文化资源与旅游产品之间是否存在转化不足或文化支撑不足的现象。换言之，本节结果更适合被理解为文旅融合潜力的初步筛查和错位诊断，而不是对文旅融合成熟度的最终评价。

\subsection{文化记忆指数（CMI）}\label{ux6587ux5316ux8bb0ux5fc6ux6307ux6570cmi}

文化记忆指数用于表达某一文化载体在典籍文本和知识图谱中的可见程度。前一阶段仅使用实体提及频次，能够反映``被记载多少''，但没有体现该实体在图谱关系中的连接位置。因此，本研究将 CMI 拆分为两个部分：一是频次厚度，二是中心度。前者反映载体相关实体在典籍中出现的强度，后者反映这些实体在知识图谱中直接相连的节点数量（中心度代表与它相连的节点多少）。

对每一条载体 c，按名称匹配得到知识图谱实体集合 Ec。设实体 e 的提及频次为 Me，则载体 c 的原始频次厚度为：

\[
Th(c) = \sum_{e \in E_{c}}^{}M_{e}
\]

为降低少数高频实体对结果的拉动，本文先取 \(\log(1 + Th(c))\)，再按 Min-Max 方法归一到 {[}0, 100{]}，得到频次分值 Fc：

\[
F_{c} = \frac{\log(1 + Th(c)) - \log(1 + Th_{\min})}{\log(1 + Th_{\max}) - \log(1 + Th_{\min})} \times 100
\]

同时，设 N(Ec) 为与 Ec 中实体直接相连的去重邻接节点集合，则载体 c 的中心度为：

\[
d_{c} = |N(E_{c})|
\]

同样对 dc 取 log(1 + dc) 后归一化，得到中心度分值 Dc。最终 CMI 采用 0.6 和 0.4 的阶段性权重合成：

\[
CMI_{c} = 0.6F_{c} + 0.4D_{c}
\]

这样处理后，CMI 同时纳入了``典籍中出现多少''和``图谱中连接多少''两个方面。专家打分在本轮不作为实测变量纳入，主要原因是目前没有形成可追溯的专家样本。如果直接加入专家分值，反而会削弱结果的可复核性。因此，专家打分更适合作为后续权重校准和敏感性检验接口，用于验证 0.6 / 0.4 的权重设定是否需要调整。

\begin{figure}
\centering
\includegraphics[width=5.90278in,height=3.02042in]{ms_thesis/image6.png}
\caption{文化记忆指数（CMI）Top 15：频次厚度与中心度}
\end{figure}

图 5.2 显示，CMI 高值载体主要集中在康有为故居、西樵山相关节点、九江镇烟桥村以及简村等传统文化空间。这些载体不仅在典籍中有较高提及频次，而且在知识图谱中能够连接到较多人物、地点、事件和文献实体，说明其文化记忆并不是单一文本频次造成的，而是具有一定关系网络支撑。

\subsection{官方认证指数（OAI）}\label{ux5b98ux65b9ux8ba4ux8bc1ux6307ux6570oai}

官方认证指数用于表达文化载体在保护名录、非遗名录或历史文化资源名录中的制度性位置。赋分规则如下：

全国重点文物保护单位、国家级非遗、中国历史文化名镇赋值为 100；

省级文物保护单位、省级非遗、广东省历史文化名村和省级传统村落赋值为 75；

市级文物保护单位和市级非遗赋值为 50；

区级认定资源赋值为 25；

未进入上述官方认定体系的资源赋值为 0。

同一载体若同时具备多重认证，取其最高等级作为 OAI 值。这样处理可以避免重复计分，同时保留最高保护层级对资源识别的影响。

\begin{figure}
\centering
\includegraphics[width=5.2in,height=3.68913in]{ms_thesis/image7.png}
\caption{官方认证指数（OAI）分值分布}
\end{figure}

图 5.3 显示，165 条核心样本中，OAI 为 50 的载体数量最多，共 73 条；OAI 为 75 的载体有 34 条，OAI 为 100 的载体有 6 条；OAI 为 25 的载体有 15 条。另有 37 条载体没有进入本研究所采用的分级认证体系，OAI 记为 0。这说明官方认证指数能够提供制度性识别依据，但并不能覆盖所有地方文化空间，仍需要与 CMI 和 THI 共同判断。

\subsection{旅游热度指数（THI）}\label{ux65c5ux6e38ux70edux5ea6ux6307ux6570thi}

旅游热度指数由 500 m 缓冲范围内的 POI 数量 \(N_{c}\)、相关 POI 平均评分 \(R_{c}\)和评论总量 \(V_{c}\)共同构成。三项指标分别归一到 {[}0, 100{]}，其中 POI 数量和评论总量先取 \(\log(1 + x)\)，以削弱极端高值点的影响。综合指数按 40\%、20\%、40\% 加权：

\[
THI_{c} = 0.4 \times \widehat{\log(1 + N_{c})} + 0.2 \times \widehat{R_{c}} + 0.4 \times \widehat{\log(1 + V_{c})}
\]

\begin{figure}
\centering
\includegraphics[width=5.90278in,height=4.24002in]{ms_thesis/image8.png}
\caption{旅游热度指数（THI）Top 15}
\end{figure}

500 m 缓冲半径与后续网格分析保持一致，大致对应步行可达范围内的周边旅游供给。评分和评论量来自平台数据，能够反映游客侧的可见热度，但也会受到平台覆盖、评论习惯和 POI 拆分方式的影响。更重要的是，本文在主体指标体系中尚未对评论文本中的文化关键词、地方叙事和游客文化感知进行语义识别，因此 THI 只能表示旅游热度，不宜直接等同于文化体验质量或文旅融合深度。

\subsection{错位指数（MI）}\label{ux9519ux4f4dux6307ux6570mi}

错位指数用于比较旅游热度与文化家底之间的相对差异，其形式为：

\[
MI_{c} = THI_{c} - \alpha \cdot CMI_{c} - \beta \cdot OAI_{c},\quad\alpha = \beta = 0.5
\]

\begin{figure}
\centering
\includegraphics[width=5.5in,height=3.77354in]{ms_thesis/image9.png}
\caption{文化---旅游错位指数（MI）分布}
\end{figure}

MI 为正，表示旅游热度相对强于文化记忆和官方认证，可能出现``旅游强、文化解释弱''的空心景点倾向；MI 为负，表示文化记忆和官方认证较强，但尚未转化为相应旅游热度，可视为沉睡潜力倾向。需要强调的是，错位并不自动等同于可干预潜力。部分负向错位可能来自交通可达性、空间承载条件或市场规模等基础约束，只有在具备展示转化、线路组织、运营活动和政策支持等条件时，才更适合被判断为可释放的文旅融合潜力。本文采用 alpha = beta = 0.5 的等权设定，是为了在文化记忆和官方认证之间保持基本平衡。第 6.3 节将进一步通过相关分析讨论该设定的解释边界，第 6.6 节则通过权重敏感性检验观察分类结果的稳定性。

\subsection{结果分层}\label{ux7ed3ux679cux5206ux5c42}

根据 CMI、THI 与 MI 的组合关系，本文将 165 条文化载体划分为四类（表 5.2）。其中，核心耦合区强调文化记忆与旅游热度的双高重合，但不等同于文旅融合已经成熟，更准确地说是具有较高优先深化价值的潜力区；沉睡潜力区强调文化资源较强但旅游转化不足；空心景点区强调旅游热度较高但文化支撑相对不足；其余样本归入一般耦合区，用作过渡性类型。

\msttableentry{5.2}{错位分层四象限的载体计数}
\textbf{表 5.2　错位分层四象限的载体计数}

{\def\LTcaptype{none} % do not increment counter
\begin{longtable}[]{@{}
  >{\centering\arraybackslash}p{(\linewidth - 6\tabcolsep) * \real{0.1538}}
  >{\centering\arraybackslash}p{(\linewidth - 6\tabcolsep) * \real{0.5641}}
  >{\centering\arraybackslash}p{(\linewidth - 6\tabcolsep) * \real{0.1538}}
  >{\centering\arraybackslash}p{(\linewidth - 6\tabcolsep) * \real{0.1282}}@{}}
\toprule\noalign{}
\endhead
\bottomrule\noalign{}
\endlastfoot
\textbf{类别} & \textbf{判别规则} & \textbf{载体数} & \textbf{占比} \\
核心耦合区 & MI 绝对值 \textless= 10，且 CMI \textgreater= 50、THI \textgreater= 50 & 5 & 3.0 \% \\
一般耦合区 & 未达到核心耦合、沉睡潜力或空心景点判别阈值 & 57 & 34.5 \% \\
沉睡潜力区 & MI \textless{} −10，且 CMI \textgreater{} THI & 54 & 32.7 \% \\
空心景点区 & MI \textgreater{} 10，且 THI \textgreater{} CMI & 49 & 29.7 \% \\
合计 & --- & 165 & 100 \% \\
\end{longtable}
}

165 条载体的 MI 最小值为 −85.72，第一四分位数为 −23.40，中位数为 −3.59，第三四分位数为 16.66，最大值为 75.37。整体上，沉睡潜力区和空心景点区合计 103 条，占 62.4\%，说明南海区文化资源与旅游产品之间存在较明显的结构性不均衡。

\begin{figure}
\centering
\includegraphics[width=5.90551in,height=4.16209in]{ms_thesis/image10.png}
\caption{500 m 网格文化---旅游分布与错位分层}
\end{figure}

图 5.6 中，横轴表示文化厚度，纵轴表示旅游热度。黄色星标标出 1 跳知识图谱口径下的核心耦合网格，用于区分一般高值点与已经形成``文化厚度---旅游热度''双高重叠的重点节点。该图的价值不在于替代空间分布图，而在于把不同错位类型放到同一坐标系中进行比较。需要注意的是，双高重叠只说明文化基础和旅游热度同时较强，后续仍需通过评论语义、产品内容和现场体验判断其文化转化深度。

\subsection{典型案例}\label{ux5178ux578bux6848ux4f8b}

为了避免仅停留在类型数量层面，本文进一步列出三类典型样本。表 5.3 至表 5.5 均按 MI 的方向和绝对值选取代表案例，用于说明不同类型的形成机制。

\msttableentry{5.3}{沉睡潜力区典型载体}
\textbf{表 5.3　沉睡潜力区典型载体}

{\def\LTcaptype{none} % do not increment counter
\begin{longtable}[]{@{}
  >{\centering\arraybackslash}p{(\linewidth - 10\tabcolsep) * \real{0.2687}}
  >{\centering\arraybackslash}p{(\linewidth - 10\tabcolsep) * \real{0.2090}}
  >{\centering\arraybackslash}p{(\linewidth - 10\tabcolsep) * \real{0.1343}}
  >{\centering\arraybackslash}p{(\linewidth - 10\tabcolsep) * \real{0.1343}}
  >{\centering\arraybackslash}p{(\linewidth - 10\tabcolsep) * \real{0.1343}}
  >{\centering\arraybackslash}p{(\linewidth - 10\tabcolsep) * \real{0.1194}}@{}}
\toprule\noalign{}
\endhead
\bottomrule\noalign{}
\endlastfoot
\textbf{载体名称} & \textbf{所属片区} & \textbf{CMI} & \textbf{OAI} & \textbf{THI} & \textbf{MI} \\
西樵镇松塘村 & 西樵 & 80.39 & 100.0 & 4.48 & −85.72 \\
九江镇烟桥烟南村 & 九江 & 84.51 & 100.0 & 7.09 & −85.16 \\
丹灶镇仙岗社区 & 丹灶 & 69.35 & 100.0 & 4.48 & −80.20 \\
赛龙舟（九江传统龙舟） & 九江 & 78.11 & 75.0 & 0.00 & −76.55 \\
西樵山仙姑亭 & 西樵 & 98.55 & 50.0 & 0.00 & −74.28 \\
\end{longtable}
}

沉睡潜力区的共同特征是文化记忆和官方认证较强，但缓冲范围内 POI、评分和评论量较低。以松塘村、烟桥村等历史村落为例，其文化价值并不低，问题主要在于旅游产品组织、到达条件和可体验内容尚未充分转化。

\msttableentry{5.4}{空心景点区典型载体}
\textbf{表 5.4　空心景点区典型载体}

{\def\LTcaptype{none} % do not increment counter
\begin{longtable}[]{@{}cccccc@{}}
\toprule\noalign{}
\endhead
\bottomrule\noalign{}
\endlastfoot
\textbf{载体名称} & \textbf{所属片区} & \textbf{CMI} & \textbf{OAI} & \textbf{THI} & \textbf{MI} \\
南海农谚 & 桂城街道 & 0.00 & 25.0 & 87.87 & 75.37 \\
南海广式旺阁酱油酿造技艺 & 桂城街道 & 0.00 & 25.0 & 70.00 & 57.50 \\
大伸市 & 九江 & 0.00 & 0.0 & 54.82 & 54.82 \\
光分亭 & 西樵 & 0.00 & 50.0 & 77.57 & 52.57 \\
红棉公园民国建筑群 & 九江 & 0.00 & 0.0 & 51.30 & 51.30 \\
\end{longtable}
}

空心景点区并不等同于``没有文化价值''。它更准确地说明：旅游热度已经高于当前可识别的文化记忆和官方认证支撑。部分样本的 CMI 接近 0，说明其文化叙事尚未在典籍文本或知识图谱中形成较强可见度；也有样本如南海农谚，虽具备区级非遗属性，但在本轮典籍频次和图谱直接连接度中识别不足，仍表现出文化解释与旅游消费之间的不均衡。

\msttableentry{5.5}{核心耦合区典型载体}
\textbf{表 5.5　核心耦合区典型载体}

{\def\LTcaptype{none} % do not increment counter
\begin{longtable}[]{@{}
  >{\centering\arraybackslash}p{(\linewidth - 10\tabcolsep) * \real{0.2687}}
  >{\centering\arraybackslash}p{(\linewidth - 10\tabcolsep) * \real{0.2090}}
  >{\centering\arraybackslash}p{(\linewidth - 10\tabcolsep) * \real{0.1343}}
  >{\centering\arraybackslash}p{(\linewidth - 10\tabcolsep) * \real{0.1343}}
  >{\centering\arraybackslash}p{(\linewidth - 10\tabcolsep) * \real{0.1343}}
  >{\centering\arraybackslash}p{(\linewidth - 10\tabcolsep) * \real{0.1194}}@{}}
\toprule\noalign{}
\endhead
\bottomrule\noalign{}
\endlastfoot
\textbf{载体名称} & \textbf{所属片区} & \textbf{CMI} & \textbf{OAI} & \textbf{THI} & \textbf{MI} \\
云泉仙馆 & 西樵 & 70.25 & 75.0 & 77.78 & 5.15 \\
西樵山百步云梯 & 西樵 & 98.58 & 50.0 & 66.51 & −7.78 \\
西樵山抗日阵亡将士暨死难同胞纪念碑 & 西樵 & 98.61 & 50.0 & 65.58 & −8.73 \\
九江吴家大院 & 九江 & 69.28 & 75.0 & 66.04 & −6.10 \\
平洲传统玉器制作技艺 & 桂城街道 & 57.28 & 50.0 & 56.63 & 2.99 \\
\end{longtable}
}

核心耦合区的样本数量较少，仅占 3.0\%。这些载体的文化记忆、官方认证和旅游热度较为接近，说明其同时具备较强文化基础和一定游客热度，是后续文旅融合深化的优先对象。但这一结果并不意味着其文化内容已经被游客充分识别，也不意味着融合已经成熟。它们更适合作为产品组织、线路设计和评论语义复核的重点节点，不宜简单复制其模式，因为不同镇街的交通条件、周边业态和文化载体类型存在差异。

\subsection{层级层面观察}\label{ux5c42ux7ea7ux5c42ux9762ux89c2ux5bdf}

从层级和类型看，12 条历史文化名村与传统村落样本全部落入沉睡潜力区，MI 中位值为 −71.82，说明村落类文化资源在文化识别和官方认证上较强，但旅游热度转化仍明显不足。非遗项目内部也呈现差异：36 项具备空间锚点的非遗中，18 项落入沉睡潜力区，省级非遗的 MI 中位值为 −28.85，市级非遗为 −17.50，区级非遗为 −5.21，说明等级越高并不必然对应更高旅游热度。不可移动文物中，一般耦合区 31 处、空心景点区 28 处、沉睡潜力区 17 处、核心耦合区 4 处，市级文物在空心景点区数量较多，提示部分文物周边已有较强旅游活动，但文化阐释和保护展示之间仍需进一步衔接。由此可见，官方等级应作为筛选和解释变量，而不宜被直接等同于旅游开发优先级。

\section{景区等级与游客体验的关联分析}\label{ux666fux533aux7b49ux7ea7ux4e0eux6e38ux5ba2ux4f53ux9a8cux7684ux5173ux8054ux5206ux6790}

\subsection{数据说明}\label{ux6570ux636eux8bf4ux660e}

为进一步观察旅游品牌等级与游客体验之间的关系，本文选取南海区 16 处已知 A 级旅游景区作为校准样本，其中包括 1 处 5A、5 处 4A、8 处 3A 和 2 处 2A。由于 A 级景区标签并不直接包含在 POI 原始分类中，本文依据公开资料整理景区名单，并将其与平台 POI 进行名称和位置匹配。需要注意的是，部分景区在平台中存在拆分点位、评分缺失或评论量偏低等情况，因此本节结果只用于提示性分析，不作强因果判断。

\subsection{相关性观察}\label{ux76f8ux5173ux6027ux89c2ux5bdf}

\msttableentry{5.6}{A 级景区相关系数（n = 16）}
\textbf{表 5.6　A 级景区相关系数（n = 16）}

{\def\LTcaptype{none} % do not increment counter
\begin{longtable}[]{@{}ccc@{}}
\toprule\noalign{}
\endhead
\bottomrule\noalign{}
\endlastfoot
\textbf{指标对} & \textbf{Pearson} & \textbf{Spearman} \\
level\_score \textasciitilde{} rating & 0.319 & 0.522 \\
level\_score \textasciitilde{} review\_count & 0.546 & 0.554 \\
level\_score \textasciitilde{} positive\_rate & 0.206 & 0.230 \\
rating \textasciitilde{} review\_count & 0.270 & 0.392 \\
\end{longtable}
}

表 5.6 显示，景区等级分数与评论量存在中等正相关（Pearson r = 0.546，Spearman r = 0.554），与平台评分的 Spearman 相关系数为 0.522；与正向率的相关性则较弱。由于样本量仅为 16，且部分平台评价存在缺失，上述结果只能说明评级等级与评论关注度、游客评分之间存在可观察联系，不能直接证明评级会导致游客满意度提高。

\begin{figure}
\centering
\includegraphics[width=4.83333in,height=2.75208in]{ms_thesis/image11.png}
\caption{景区评级与评论热度、评分的关系}
\end{figure}

\subsection{对规划的启示}\label{ux5bf9ux89c4ux5212ux7684ux542fux793a}

从规划角度看，A 级景区申报可以作为文旅转化的一种制度工具，尤其适用于文化记忆较强、空间载体明确、但品牌识别不足的资源。但评级本身并不能替代内容建设。对于沉睡潜力区，应优先完善交通接驳、游线组织和解说系统；对于空心景点区，则需要补足地方文化叙事和体验内容。也就是说，评级可以提高资源的可见度，但游客体验的改善仍取决于具体产品是否能够把文化内容转化为可感知、可停留、可传播的游览过程。
